\section{System analysis}

This chapter presents a systematic analysis of the problem domain, following established system analysis principles.

\subsection{System representation}

The system under study can be represented as a ``black box'' with clearly defined inputs and outputs.
The inputs include raw data, configuration parameters, and user requirements.
The outputs include processed results, performance metrics, and generated reports.

\subsection{System functions}

The system performs the following primary functions:
\begin{itemize}
  \item data ingestion and preprocessing;
  \item analysis and transformation;
  \item result generation and visualization;
  \item logging and monitoring;
  \item reporting and export.
\end{itemize}

\subsection{System decomposition}

The system is decomposed into four subsystems:
\begin{itemize}
  \item data processing subsystem;
  \item analysis engine subsystem;
  \item user interface subsystem;
  \item monitoring and logging subsystem.
\end{itemize}

\begin{figure}[H]
    \centering
    \begin{tikzpicture}[
        box/.style={rectangle, draw, minimum width=3cm, minimum height=0.8cm, align=center},
        >=Stealth
    ]
        \node[box] (sys) {System};
        \node[box, below left=1cm and 0.5cm of sys] (dp) {Data Processing};
        \node[box, below=1cm of sys] (ae) {Analysis Engine};
        \node[box, below right=1cm and 0.5cm of sys] (ui) {User Interface};
        \node[box, below=1cm of ae] (ml) {Monitoring \& Logging};
        \draw[->] (sys) -- (dp);
        \draw[->] (sys) -- (ae);
        \draw[->] (sys) -- (ui);
        \draw[->] (ae) -- (ml);
    \end{tikzpicture}
    \caption{System decomposition diagram} \label{fig:decomposition}
\end{figure}

\subsection{Hierarchical structure}

The hierarchical organization follows the principle of separation of concerns.
At the top level, the orchestrator manages the overall workflow.
The middle layer contains domain-specific processing modules.
The bottom layer handles infrastructure concerns such as data storage and network communication.

\begin{figure}[H]
    \centering
    \begin{tikzpicture}[
        box/.style={rectangle, draw, minimum width=2.5cm, minimum height=0.7cm, align=center},
        >=Stealth
    ]
        \node[box] (orch) {Orchestrator};
        \node[box, below left=0.8cm and 1cm of orch] (proc) {Processing};
        \node[box, below right=0.8cm and 1cm of orch] (report) {Reporting};
        \node[box, below left=0.8cm and 0cm of proc] (store) {Storage};
        \node[box, below right=0.8cm and 0cm of proc] (net) {Network};
        \draw[->] (orch) -- (proc);
        \draw[->] (orch) -- (report);
        \draw[->] (proc) -- (store);
        \draw[->] (proc) -- (net);
    \end{tikzpicture}
    \caption{System hierarchy diagram} \label{fig:hierarchy}
\end{figure}

\subsection{Summary}

The system analysis reveals that a modular architecture with clear boundaries between subsystems provides the best foundation for achieving the research objectives.

\clearpage
